\section{Controlled-mechanism results}
\label{app:controlled}

The controlled studies isolate the mechanism behind the real-data negatives in settings where
the hidden state and the optimal action are known, so a null cannot be blamed on market noise.

\paragraph{Action-relevant ambiguity (E-210).} Table~\ref{tab:e210} varies the training size
and the label family. Restricted summaries incur a decision regret (for example $0.075$ at the
smallest universal setting) that shrinks with data, while the matched direct and
posterior-restoration coordinates are equal, so no restoration-specific advantage appears. The
latent-mean reconstruction MSE moves independently of the decision cost: recovering the latent
mean better does not lower the decision cost, which is the controlled analogue of the Q16
market result.

\begin{table}[h]\centering
\caption{Controlled action-ambiguity (E-210): ranked probability score, primary decision cost,
decision regret, and latent-mean reconstruction MSE by training size and family. Restricted
summaries incur regret while matched coordinates are equal. Source:
\texttt{experiments/results/e210\_finite.csv}.}
\label{tab:e210}
\begin{tabular}{llrrrr}
\toprule
$n_{\mathrm{train}}$ & family & RPS & primary cost & regret & latent MSE \\
\midrule
32 & universal & 0.201384 & 0.435500 & 0.010500 & 1.298229 \\
32 & symmetric & 0.194937 & 0.425750 & 0.000750 & 1.250000 \\
128 & universal & 0.196002 & 0.425656 & 0.000656 & 1.267059 \\
128 & symmetric & 0.193729 & 0.425000 & 0.000000 & 1.250000 \\
512 & universal & 0.194103 & 0.425000 & 0.000000 & 1.255075 \\
512 & symmetric & 0.193440 & 0.425000 & 0.000000 & 1.250000 \\
\bottomrule
\end{tabular}\end{table}


\paragraph{Point versus distributional fill prediction (E-258).} Table~\ref{tab:e258} compares a
direct predictor, a point-reconstruction predictor, and a distributional-reconstruction
predictor across conditions and training sizes. Hidden-queue uncertainty matters, the
point-reconstruction predictor is clearly worse, but the direct and distributional predictors
coincide at $512$ draws and the direct-compiled predictor matches the distributional one
exactly, so reconstruction carries no isolated advantage over predicting the decision directly.

\begin{longtable}{lllrrrr}
\caption{Controlled fill-prediction (E-258), selection role: Brier, log loss, calibration
error, and rank MSE by condition, training size, and predictor. Direct and distributional
predictors coincide at 512 draws; reconstruction shows no isolated advantage. Source:
\texttt{experiments/results/e258\_metrics.csv}.}
\label{tab:e258}\\
\toprule
condition & size & model & Brier & log loss & ECE & rank MSE \\
\midrule
\endfirsthead
\toprule
condition & size & model & Brier & log loss & ECE & rank MSE \\
\midrule
\endhead
\bottomrule
\endfoot
informative & 128 & unconditional\_y & 0.236403 & 0.665809 & 0.063582 & -- \\
informative & 128 & direct\_y & 0.218317 & 0.627603 & 0.063981 & -- \\
informative & 128 & point\_y & 0.272500 & 0.756857 & 0.205469 & 0.562500 \\
informative & 128 & distribution\_y & 0.218936 & 0.628852 & 0.065705 & 0.557968 \\
informative & 128 & direct\_compiled\_y & 0.218936 & 0.628852 & 0.065705 & -- \\
informative & 128 & current\_rich\_y & 0.145427 & 0.468343 & 0.135528 & -- \\
informative & 128 & distribution\_h & 0.213186 & 0.613884 & 0.060247 & 0.548696 \\
informative & 128 & point\_mean\_h & 0.272500 & 0.756857 & 0.205469 & 0.562500 \\
informative & 128 & direct\_soft\_h & 0.213186 & 0.613884 & 0.060247 & -- \\
informative & 128 & known\_mean\_point & 0.272500 & 0.756857 & 0.205469 & 0.562500 \\
informative & 128 & exact\_coarse\_current & 0.216287 & 0.620818 & 0.073906 & -- \\
informative & 128 & exact\_coarse\_source & 0.216287 & 0.620818 & 0.073906 & -- \\
informative & 128 & exact\_rich & 0.128750 & 0.420450 & 0.012500 & -- \\
informative & 128 & context\_only & 0.233359 & 0.658356 & 0.085156 & -- \\
informative & 128 & exact\_unconditional & 0.239219 & 0.671521 & 0.082813 & -- \\
informative & 512 & unconditional\_y & 0.237339 & 0.667713 & 0.070556 & -- \\
informative & 512 & direct\_y & 0.210430 & 0.608710 & 0.054102 & -- \\
informative & 512 & point\_y & 0.272500 & 0.756857 & 0.205469 & 0.562500 \\
informative & 512 & distribution\_y & 0.210430 & 0.608710 & 0.054102 & 0.546351 \\
informative & 512 & direct\_compiled\_y & 0.210430 & 0.608710 & 0.054102 & -- \\
informative & 512 & current\_rich\_y & 0.129540 & 0.422807 & 0.041310 & -- \\
informative & 512 & distribution\_h & 0.214822 & 0.617609 & 0.076785 & 0.550128 \\
informative & 512 & point\_mean\_h & 0.272500 & 0.756857 & 0.205469 & 0.562500 \\
informative & 512 & direct\_soft\_h & 0.214822 & 0.617609 & 0.076785 & -- \\
informative & 512 & known\_mean\_point & 0.272500 & 0.756857 & 0.205469 & 0.562500 \\
informative & 512 & exact\_coarse\_current & 0.216287 & 0.620818 & 0.073906 & -- \\
informative & 512 & exact\_coarse\_source & 0.216287 & 0.620818 & 0.073906 & -- \\
informative & 512 & exact\_rich & 0.128750 & 0.420450 & 0.012500 & -- \\
informative & 512 & context\_only & 0.233359 & 0.658356 & 0.085156 & -- \\
informative & 512 & exact\_unconditional & 0.239219 & 0.671521 & 0.082813 & -- \\
zero & 128 & unconditional\_y & 0.260971 & 0.716023 & 0.107332 & -- \\
zero & 128 & direct\_y & 0.233780 & 0.660119 & 0.099132 & -- \\
zero & 128 & point\_y & 0.227187 & 0.658831 & 0.111719 & 0.460938 \\
zero & 128 & distribution\_y & 0.233780 & 0.660119 & 0.099132 & 0.540339 \\
zero & 128 & direct\_compiled\_y & 0.233780 & 0.660119 & 0.099132 & -- \\
zero & 128 & current\_rich\_y & 0.141311 & 0.449890 & 0.092970 & -- \\
zero & 128 & distribution\_h & 0.219565 & 0.630344 & 0.074597 & 0.485886 \\
zero & 128 & point\_mean\_h & 0.227187 & 0.658831 & 0.111719 & 0.460938 \\
zero & 128 & direct\_soft\_h & 0.219565 & 0.630344 & 0.074597 & -- \\
zero & 128 & known\_mean\_point & 0.227187 & 0.658831 & 0.111719 & 0.460938 \\
zero & 128 & exact\_coarse\_current & 0.216172 & 0.623795 & 0.038281 & -- \\
zero & 128 & exact\_coarse\_source & 0.216172 & 0.623795 & 0.038281 & -- \\
zero & 128 & exact\_rich & 0.125098 & 0.411581 & 0.034766 & -- \\
zero & 128 & context\_only & 0.216172 & 0.623795 & 0.038281 & -- \\
zero & 128 & exact\_unconditional & 0.250156 & 0.693469 & 0.026562 & -- \\
zero & 512 & unconditional\_y & 0.251807 & 0.696824 & 0.048547 & -- \\
zero & 512 & direct\_y & 0.219139 & 0.629516 & 0.062349 & -- \\
zero & 512 & point\_y & 0.227187 & 0.658831 & 0.111719 & 0.460938 \\
zero & 512 & distribution\_y & 0.219139 & 0.629516 & 0.062349 & 0.475028 \\
zero & 512 & direct\_compiled\_y & 0.219139 & 0.629516 & 0.062349 & -- \\
zero & 512 & current\_rich\_y & 0.130753 & 0.428611 & 0.098630 & -- \\
zero & 512 & distribution\_h & 0.216958 & 0.625334 & 0.041112 & 0.466307 \\
zero & 512 & point\_mean\_h & 0.227187 & 0.658831 & 0.111719 & 0.460938 \\
zero & 512 & direct\_soft\_h & 0.216958 & 0.625334 & 0.041112 & -- \\
zero & 512 & known\_mean\_point & 0.227187 & 0.658831 & 0.111719 & 0.460938 \\
zero & 512 & exact\_coarse\_current & 0.216172 & 0.623795 & 0.038281 & -- \\
zero & 512 & exact\_coarse\_source & 0.216172 & 0.623795 & 0.038281 & -- \\
zero & 512 & exact\_rich & 0.125098 & 0.411581 & 0.034766 & -- \\
zero & 512 & context\_only & 0.216172 & 0.623795 & 0.038281 & -- \\
zero & 512 & exact\_unconditional & 0.250156 & 0.693469 & 0.026562 & -- \\
\end{longtable}


\paragraph{Fixed-model persistence (E-271).} Table~\ref{tab:e271} reports the persistence check
on one admitted BTC day. The strict joint condition is contradicted: raw and favorable-class
contrasts are positive while the fixed-class contrast is $-0.01337$, and a uniform predictor
beats every learned breakout arm on observed-event loss. One source-exposed day cannot support a
population claim, but it is consistent with the reconstruction and sufficiency negatives: the
richer, learned arm does not persist as the better decision.

\begin{longtable}{llllrlr}
\caption{Fixed-model persistence (E-271) on BTC 2026-02-01: per-endpoint values by family,
kind, and metric. The fixed-class contrast is negative while raw and favorable-class contrasts
are positive. Source: \texttt{experiments/results/e271\_main.csv}.}
\label{tab:e271}\\
\toprule
family & kind & metric & endpoint & value & status & $n$ \\
\midrule
\endfirsthead
\toprule
family & kind & metric & endpoint & value & status & $n$ \\
\midrule
\endhead
\bottomrule
\endfoot
breakout & method & C & raw & 0.377905 & evaluable & 880 \\
breakout & method & C & frozen\_mixture & 0.303310 & evaluable & 880 \\
breakout & method & C & F & 0.277484 & evaluable & 880 \\
breakout & method & C & A & 0.618632 & evaluable & 880 \\
breakout & method & C & N & 0.243273 & evaluable & 880 \\
breakout & method & P & raw & 0.376674 & evaluable & 880 \\
breakout & method & P & frozen\_mixture & 0.295825 & evaluable & 880 \\
breakout & method & P & F & 0.293166 & evaluable & 880 \\
breakout & method & P & A & 0.610970 & evaluable & 880 \\
breakout & method & P & N & 0.220948 & evaluable & 880 \\
breakout & method & R & raw & 0.378487 & evaluable & 880 \\
breakout & method & R & frozen\_mixture & 0.289943 & evaluable & 880 \\
breakout & method & R & F & 0.294158 & evaluable & 880 \\
breakout & method & R & A & 0.627595 & evaluable & 880 \\
breakout & method & R & N & 0.205180 & evaluable & 880 \\
breakout & method & onehot-F & raw & 0.568182 & evaluable & 880 \\
breakout & method & onehot-F & frozen\_mixture & 0.659259 & evaluable & 880 \\
breakout & method & onehot-F & F & 0.000000 & evaluable & 880 \\
breakout & method & onehot-F & A & 1.000000 & evaluable & 880 \\
breakout & method & onehot-F & N & 1.000000 & evaluable & 880 \\
breakout & method & onehot-A & raw & 0.680682 & evaluable & 880 \\
breakout & method & onehot-A & frozen\_mixture & 0.871111 & evaluable & 880 \\
breakout & method & onehot-A & F & 1.000000 & evaluable & 880 \\
breakout & method & onehot-A & A & 0.000000 & evaluable & 880 \\
breakout & method & onehot-A & N & 1.000000 & evaluable & 880 \\
breakout & method & onehot-N & raw & 0.751136 & evaluable & 880 \\
breakout & method & onehot-N & frozen\_mixture & 0.469630 & evaluable & 880 \\
breakout & method & onehot-N & F & 1.000000 & evaluable & 880 \\
breakout & method & onehot-N & A & 1.000000 & evaluable & 880 \\
breakout & method & onehot-N & N & 0.000000 & evaluable & 880 \\
breakout & method & uniform & raw & 0.333333 & evaluable & 880 \\
breakout & method & uniform & frozen\_mixture & 0.333333 & evaluable & 880 \\
breakout & method & uniform & F & 0.333333 & evaluable & 880 \\
breakout & method & uniform & A & 0.333333 & evaluable & 880 \\
breakout & method & uniform & N & 0.333333 & evaluable & 880 \\
breakout & method & train-frequency & raw & 0.386720 & evaluable & 880 \\
breakout & method & train-frequency & frozen\_mixture & 0.292995 & evaluable & 880 \\
breakout & method & train-frequency & F & 0.366264 & evaluable & 880 \\
breakout & method & train-frequency & A & 0.578116 & evaluable & 880 \\
breakout & method & train-frequency & N & 0.176634 & evaluable & 880 \\
breakout & contrast & P-C & raw & -0.001231 & evaluable & 880 \\
breakout & contrast & P-C & frozen\_mixture & -0.007485 & evaluable & 880 \\
breakout & contrast & P-C & F & 0.015682 & evaluable & 880 \\
breakout & contrast & P-C & A & -0.007662 & evaluable & 880 \\
breakout & contrast & P-C & N & -0.022325 & evaluable & 880 \\
breakout & contrast & R-C & raw & 0.000582 & evaluable & 880 \\
breakout & contrast & R-C & frozen\_mixture & -0.013367 & evaluable & 880 \\
breakout & contrast & R-C & F & 0.016674 & evaluable & 880 \\
breakout & contrast & R-C & A & 0.008963 & evaluable & 880 \\
breakout & contrast & R-C & N & -0.038093 & evaluable & 880 \\
breakout & contrast & R-P & raw & 0.001813 & evaluable & 880 \\
breakout & contrast & R-P & frozen\_mixture & -0.005882 & evaluable & 880 \\
breakout & contrast & R-P & F & 0.000992 & evaluable & 880 \\
breakout & contrast & R-P & A & 0.016625 & evaluable & 880 \\
breakout & contrast & R-P & N & -0.015768 & evaluable & 880 \\
breakout & contrast & C-frequency & raw & -0.008815 & evaluable & 880 \\
breakout & contrast & C-frequency & frozen\_mixture & 0.010314 & evaluable & 880 \\
breakout & contrast & C-frequency & F & -0.088780 & evaluable & 880 \\
breakout & contrast & C-frequency & A & 0.040516 & evaluable & 880 \\
breakout & contrast & C-frequency & N & 0.066639 & evaluable & 880 \\
breakout & contrast & P-frequency & raw & -0.010046 & evaluable & 880 \\
breakout & contrast & P-frequency & frozen\_mixture & 0.002830 & evaluable & 880 \\
breakout & contrast & P-frequency & F & -0.073098 & evaluable & 880 \\
breakout & contrast & P-frequency & A & 0.032854 & evaluable & 880 \\
breakout & contrast & P-frequency & N & 0.044314 & evaluable & 880 \\
breakout & contrast & R-frequency & raw & -0.008233 & evaluable & 880 \\
breakout & contrast & R-frequency & frozen\_mixture & -0.003053 & evaluable & 880 \\
breakout & contrast & R-frequency & F & -0.072106 & evaluable & 880 \\
breakout & contrast & R-frequency & A & 0.049479 & evaluable & 880 \\
breakout & contrast & R-frequency & N & 0.028546 & evaluable & 880 \\
breakout & contrast & frequency-uniform & raw & 0.053387 & evaluable & 880 \\
breakout & contrast & frequency-uniform & frozen\_mixture & -0.040338 & evaluable & 880 \\
breakout & contrast & frequency-uniform & F & 0.032931 & evaluable & 880 \\
breakout & contrast & frequency-uniform & A & 0.244782 & evaluable & 880 \\
breakout & contrast & frequency-uniform & N & -0.156699 & evaluable & 880 \\
rebound & method & C & raw & 0.339716 & evaluable & 865 \\
rebound & method & C & frozen\_mixture & 0.310577 & evaluable & 865 \\
rebound & method & C & F & 0.711610 & evaluable & 865 \\
rebound & method & C & A & 0.197967 & evaluable & 865 \\
rebound & method & C & N & 0.359078 & evaluable & 865 \\
rebound & method & P & raw & 0.341107 & evaluable & 865 \\
rebound & method & P & frozen\_mixture & 0.310627 & evaluable & 865 \\
rebound & method & P & F & 0.702873 & evaluable & 865 \\
rebound & method & P & A & 0.207901 & evaluable & 865 \\
rebound & method & P & N & 0.350677 & evaluable & 865 \\
rebound & method & R & raw & 0.366120 & evaluable & 865 \\
rebound & method & R & frozen\_mixture & 0.318122 & evaluable & 865 \\
rebound & method & R & F & 0.773648 & evaluable & 865 \\
rebound & method & R & A & 0.246041 & evaluable & 865 \\
rebound & method & R & N & 0.317585 & evaluable & 865 \\
rebound & method & onehot-F & raw & 0.809249 & evaluable & 865 \\
rebound & method & onehot-F & frozen\_mixture & 0.926298 & evaluable & 865 \\
rebound & method & onehot-F & F & 0.000000 & evaluable & 865 \\
rebound & method & onehot-F & A & 1.000000 & evaluable & 865 \\
rebound & method & onehot-F & N & 1.000000 & evaluable & 865 \\
rebound & method & onehot-A & raw & 0.462428 & evaluable & 865 \\
rebound & method & onehot-A & frozen\_mixture & 0.537688 & evaluable & 865 \\
rebound & method & onehot-A & F & 1.000000 & evaluable & 865 \\
rebound & method & onehot-A & A & 0.000000 & evaluable & 865 \\
rebound & method & onehot-A & N & 1.000000 & evaluable & 865 \\
rebound & method & onehot-N & raw & 0.728324 & evaluable & 865 \\
rebound & method & onehot-N & frozen\_mixture & 0.536013 & evaluable & 865 \\
rebound & method & onehot-N & F & 1.000000 & evaluable & 865 \\
rebound & method & onehot-N & A & 1.000000 & evaluable & 865 \\
rebound & method & onehot-N & N & 0.000000 & evaluable & 865 \\
rebound & method & uniform & raw & 0.333333 & evaluable & 865 \\
rebound & method & uniform & frozen\_mixture & 0.333333 & evaluable & 865 \\
rebound & method & uniform & F & 0.333333 & evaluable & 865 \\
rebound & method & uniform & A & 0.333333 & evaluable & 865 \\
rebound & method & uniform & N & 0.333333 & evaluable & 865 \\
rebound & method & train-frequency & raw & 0.328585 & evaluable & 865 \\
rebound & method & train-frequency & frozen\_mixture & 0.282776 & evaluable & 865 \\
rebound & method & train-frequency & F & 0.643522 & evaluable & 865 \\
rebound & method & train-frequency & A & 0.254912 & evaluable & 865 \\
rebound & method & train-frequency & N & 0.253237 & evaluable & 865 \\
rebound & contrast & P-C & raw & 0.001391 & evaluable & 865 \\
rebound & contrast & P-C & frozen\_mixture & 0.000050 & evaluable & 865 \\
rebound & contrast & P-C & F & -0.008737 & evaluable & 865 \\
rebound & contrast & P-C & A & 0.009934 & evaluable & 865 \\
rebound & contrast & P-C & N & -0.008401 & evaluable & 865 \\
rebound & contrast & R-C & raw & 0.026404 & evaluable & 865 \\
rebound & contrast & R-C & frozen\_mixture & 0.007545 & evaluable & 865 \\
rebound & contrast & R-C & F & 0.062038 & evaluable & 865 \\
rebound & contrast & R-C & A & 0.048073 & evaluable & 865 \\
rebound & contrast & R-C & N & -0.041493 & evaluable & 865 \\
rebound & contrast & R-P & raw & 0.025013 & evaluable & 865 \\
rebound & contrast & R-P & frozen\_mixture & 0.007495 & evaluable & 865 \\
rebound & contrast & R-P & F & 0.070775 & evaluable & 865 \\
rebound & contrast & R-P & A & 0.038140 & evaluable & 865 \\
rebound & contrast & R-P & N & -0.033092 & evaluable & 865 \\
rebound & contrast & C-frequency & raw & 0.011131 & evaluable & 865 \\
rebound & contrast & C-frequency & frozen\_mixture & 0.027801 & evaluable & 865 \\
rebound & contrast & C-frequency & F & 0.068088 & evaluable & 865 \\
rebound & contrast & C-frequency & A & -0.056945 & evaluable & 865 \\
rebound & contrast & C-frequency & N & 0.105841 & evaluable & 865 \\
rebound & contrast & P-frequency & raw & 0.012522 & evaluable & 865 \\
rebound & contrast & P-frequency & frozen\_mixture & 0.027851 & evaluable & 865 \\
rebound & contrast & P-frequency & F & 0.059351 & evaluable & 865 \\
rebound & contrast & P-frequency & A & -0.047011 & evaluable & 865 \\
rebound & contrast & P-frequency & N & 0.097440 & evaluable & 865 \\
rebound & contrast & R-frequency & raw & 0.037535 & evaluable & 865 \\
rebound & contrast & R-frequency & frozen\_mixture & 0.035346 & evaluable & 865 \\
rebound & contrast & R-frequency & F & 0.130126 & evaluable & 865 \\
rebound & contrast & R-frequency & A & -0.008871 & evaluable & 865 \\
rebound & contrast & R-frequency & N & 0.064348 & evaluable & 865 \\
rebound & contrast & frequency-uniform & raw & -0.004748 & evaluable & 865 \\
rebound & contrast & frequency-uniform & frozen\_mixture & -0.050557 & evaluable & 865 \\
rebound & contrast & frequency-uniform & F & 0.310189 & evaluable & 865 \\
rebound & contrast & frequency-uniform & A & -0.078421 & evaluable & 865 \\
rebound & contrast & frequency-uniform & N & -0.080096 & evaluable & 865 \\
breakout & joint\_component & R-C & raw & 0.000582 & evaluable & 880 \\
breakout & joint\_component & R-C & fixed\_class & -0.013367 & evaluable & 880 \\
breakout & joint\_component & R-C & favorable\_conditional & 0.016674 & evaluable & 880 \\
\end{longtable}

